\documentclass[11pt]{article}
\usepackage[a4paper,margin=25mm]{geometry}
\usepackage[T1]{fontenc}
\usepackage{lmodern,amsmath,amssymb,amsthm,mathtools,microtype,booktabs,enumitem}
\usepackage[colorlinks=true,linkcolor=blue,urlcolor=blue,citecolor=blue]{hyperref}
\hypersetup{pdftitle={Why zeta(3) is irrational},pdfauthor={Mathathon peer-review digest}}
\newcommand{\Z}{\mathbb Z}
\newcommand{\Q}{\mathbb Q}
\newcommand{\dd}{\,\mathrm d}
\newcommand{\ip}[2]{\langle #1,#2\rangle}
\DeclareMathOperator{\lcm}{lcm}
\newtheorem{theorem}{Theorem}
\newtheorem{lemma}{Lemma}
\newtheorem{proposition}[lemma]{Proposition}
\theoremstyle{remark}
\newtheorem*{remark}{Remark}
\setlist{itemsep=1pt,topsep=3pt}
\newcommand{\q}[1]{\medskip\noindent\textbf{#1}}
\newcommand{\forced}{\medskip\noindent\textbf{Forced or chosen?}\ }
\title{Why $\zeta(3)$ is irrational\\[2mm]
\large the proof, built one need at a time}
\author{Mathathon peer-review digest\\
\small written by Claude from the Claude and Codex submissions in \texttt{Peer Review/}}
\date{14 November 2026 (rewritten for motivation)}
\begin{document}
\maketitle

\noindent\textbf{How this document is written.} Every proof of Ap\'ery's theorem contains
objects that look as if they fell from the sky: a kernel $\frac{-\log(xy)}{1-xy}$, the polynomial
$\frac1{n!}\frac{d^n}{dx^n}[x^n(1-x)^n]$, the integers $\sum_k\binom nk^2\binom{n+k}k^2$, the
numbers $17$ and $(\sqrt2-1)^4$. Here nothing is introduced until a difficulty demands it. Each
section has the same shape: a \emph{question}; the \emph{obvious attempt} and the precise way it
fails; the \emph{construction} that the failure points to; and a note on what was \emph{forced} and
what was a choice. Whenever we reject an alternative we give the numbers.

Part I (Sections 1--8) is one complete proof, due in essence to Beukers. Part II (Sections 9--11)
shows the same objects in the language of Zudilin's paper, where they are sums instead of
integrals, and explains where Ap\'ery's recurrence comes from. No symbol is used before it is
defined. The prerequisites are integration by parts, geometric series and the AM--GM inequality.

\smallskip\noindent\textbf{Status.} The theorem and constructions are classical (Ap\'ery,
Beukers, Zudilin); nothing here is claimed to be new, and when we say an object is ``forced'' we
mean forced by the stated requirement, not that no other proof could exist. All statements were
checked by AI agents only; \emph{no human has refereed this}. Details are in Section~12.

\part*{Part I. One proof, from the beginning}

\section{What does any irrationality proof have to do?}

The number in question is $\zeta(3)=1+\frac1{2^3}+\frac1{3^3}+\cdots=1.2020569\ldots$

\q{Question.} How can one possibly prove that a number is not a fraction?

A number given by an infinite series has no visible reason to be irrational: $\sum_{m\ge1}2^{-m}=1$.
What we can use is the one rigid fact about fractions: \emph{there is a gap between $0$ and the
next integer.} Suppose $\zeta(3)=p/q$ with positive integers $p,q$. Then for any integers $A,B$,
\[
q\cdot\bigl(A\,\zeta(3)-B\bigr)=Ap-Bq\quad\text{is an integer, so it is $0$ or at least $1$ in size.}
\]
Hence $|A\zeta(3)-B|$ is either $0$ or at least $1/q$. So:

\begin{lemma}[The criterion]\label{lem:crit}
If there are integers $A_n,B_n$ with $A_n\zeta(3)-B_n\ne0$ for all $n$ and
$A_n\zeta(3)-B_n\to0$, then $\zeta(3)$ is irrational.
\end{lemma}

\forced Forced. This is not one trick among many: for \emph{any} real $\theta$, irrationality is
equivalent to the existence of such integers. (If $\theta$ is irrational, look at the $N+1$
fractional parts of $0,\theta,2\theta,\dots,N\theta$; two of them lie in the same interval of length
$1/N$, and their difference gives integers with $0<|A\theta-B|<1/N$.) So every proof, whatever it
looks like, must in the end produce integers $A_n,B_n$ making $A_n\zeta(3)-B_n$ small and
non-zero. The only question is how to find them.

Notice what is being asked. It is not that $B_n/A_n$ be close to $\zeta(3)$; that is easy. It is
that the error \emph{multiplied by the denominator $A_n$} be small. That is a much stronger demand.

\section{The obvious attempt, and the lesson it teaches}

\q{Attempt.} Cut the series off: $s_n=\sum_{m=1}^n\frac1{m^3}$. This is a fraction close to
$\zeta(3)$.

\emph{Its denominator.} To write $s_n$ as $B/A$ we need a common denominator for
$1,\frac18,\dots,\frac1{n^3}$. The cheapest is the cube of
\[
D_n=\lcm(1,2,\dots,n),\qquad\text{the smallest positive integer divisible by each of }1,\dots,n .
\]
($D_1,\dots,D_{10}=1,2,6,12,60,60,420,840,2520,2520$.) Why $D_n$ and not $n!$? Both work, but $n!$
contains each prime many times over, while $D_n$ contains each prime only to the highest power
that fits below $n$. The difference is enormous: $n!$ grows faster than any exponential, whereas
$D_n$ grows roughly like $e^n$, with $e=2.718\ldots$ (this is the prime number theorem; we shall
need, and prove, much less in Section~7).

\emph{Its error.} $\zeta(3)-s_n=\sum_{m>n}m^{-3}\ge\frac1{2(n+1)^2}$. So with $A=D_n^3$ and
$B=D_n^3s_n$:
\[
A\,\zeta(3)-B=D_n^3\bigl(\zeta(3)-s_n\bigr)\ \ge\ \frac{D_n^3}{2(n+1)^2}\ \to\ \infty .
\]
Numerically this is $3541$ at $n=5$, $7\cdot10^7$ at $n=10$ and $10^{22}$ at $n=20$.

\q{Lesson.} The error of the partial sum is only \emph{polynomially} small ($\approx n^{-2}$), but
the price of clearing denominators is \emph{exponentially} large ($D_n^3\approx e^{3n}\approx20^n$).
This sets the target for everything that follows:

\begin{quote}
\emph{Find combinations $a\,\zeta(3)-b$ whose denominators are still cleared by $D_n^3$, but which
are smaller than $20^{-n}$ --- and not zero.}
\end{quote}

\noindent So we need three things at once: (A) cheap denominators; (B) a value that is non-zero and
exponentially small; (C) a guarantee that $D_n^3$ really costs less than the smallness gained.
Partial sums have (A) and fail (B). We keep their arithmetic and look for more cancellation.

\section{A machine that produces combinations with cheap denominators}

\q{Question.} Partial sums are the only combinations $a\zeta(3)-b$ we know. How do we get a large
supply of others, all with denominator $D_n^3$, among which to search for a small one?

\q{Idea.} Write $\frac1{m^3}$ as an integral; then we can put a polynomial inside the integral,
and every polynomial gives a new combination. Integration manufactures reciprocals of integers:
\[
\int_0^1x^{m}\dd x=\frac1{m+1},\qquad\int_0^1x^{m}(-\log x)\dd x=\frac1{(m+1)^2}
\]
(the second by integration by parts). We need \emph{three} factors of $m+1$. Two variables give
two, and one logarithm gives the third:
$\iint(xy)^m(-\log x)\dd x\dd y=\frac1{(m+1)^3}$. Summing over $m\ge0$ with the geometric series
$\sum_m(xy)^m=\frac1{1-xy}$, and using $-\log x-\log y=-\log(xy)$ to make the expression symmetric
in $x$ and $y$, we arrive at the \emph{kernel}
\begin{equation}\label{kernel}
K(x,y)=\frac{-\log(xy)}{1-xy},\qquad\int_0^1\!\!\int_0^1K(x,y)\dd x\dd y=2\,\zeta(3).
\end{equation}
The factor $2$ is there because $-\log(xy)$ consists of two equal pieces, $-\log x$ and $-\log y$.
We keep it rather than carry a $\frac12$ through every formula. (All terms are positive, so
exchanging sum and integral is legitimate.)

\q{What happens when a polynomial is inserted?} It is enough to insert monomials $x^ry^s$. Write
$H_k^{(j)}=1+\frac1{2^j}+\cdots+\frac1{k^j}$ for the partial sums of $\zeta(j)$, with $H_0^{(j)}=0$.

\begin{lemma}[Moments of the kernel]\label{lem:moments}
For integers $r,s\ge0$,
\[
\int_0^1\!\!\int_0^1x^ry^sK(x,y)\dd x\dd y=
\begin{cases}
2\bigl(\zeta(3)-H_r^{(3)}\bigr)&\text{if }r=s,\\[3pt]
\dfrac{H_r^{(2)}-H_s^{(2)}}{r-s}&\text{if }r\ne s.
\end{cases}
\]
\end{lemma}
\begin{proof}
Expanding as above, the integral is $\sum_{m\ge0}\bigl[\frac1{(m+r+1)^2(m+s+1)}+
\frac1{(m+r+1)(m+s+1)^2}\bigr]$.
If $r=s$ each bracket is $\frac2{(m+r+1)^3}$: the series for $2\zeta(3)$ with its first $r$ terms
missing. If $r>s$, put $u=m+s+1$ and $d=r-s$; the bracket is
$\frac{2u+d}{u^2(u+d)^2}=\frac1d\bigl[\frac1{u^2}-\frac1{(u+d)^2}\bigr]$, and the sum over $u>s$
telescopes, leaving $\frac1d\sum_{u=s+1}^{r}\frac1{u^2}$.
\end{proof}
\noindent\emph{Why this is the heart of the arithmetic.} On the diagonal $r=s$ both variables are
shifted by the same amount, and we just get $\zeta(3)$ with a few terms removed. Off the diagonal
the two shifts differ, the sum collapses, and only a \emph{rational number} survives. So $\zeta(3)$
enters only through the diagonal, and no other irrational constant enters at all.

Now let $P(x)=\sum_{j=0}^nc_jx^j$ have integer coefficients, and define
\[
J(P)=\int_0^1\!\!\int_0^1P(x)P(y)K(x,y)\dd x\dd y .
\]
\begin{proposition}[Cheap denominators]\label{prop:arith}
$J(P)=2\bigl(\sum_jc_j^2\bigr)\zeta(3)-\beta(P)$, where $\beta(P)$ is a rational number and
$D_n^3\,\beta(P)$ is an integer.
\end{proposition}
\begin{proof}
Multiply out $P(x)P(y)$ and use Lemma~\ref{lem:moments}. The diagonal terms give
$2\sum c_j^2(\zeta(3)-H_j^{(3)})$; the denominators in $H_j^{(3)}$ are $k^3$ with $k\le n$. The
off-diagonal terms are integer multiples of $\frac1{(r-s)u^2}$ with $|r-s|\le n$ and $u\le n$.
All of these divide $D_n^3$.
\end{proof}

\q{The crucial observation.} The price $D_n^3$ depends only on the \emph{degree} of $P$, not on the
size of its coefficients. So we have $n+1$ integers $c_0,\dots,c_n$ to choose, as large as we like,
at a fixed price. For instance $P=1$ gives $J=2\zeta(3)$, and $P=1-2x$ gives
$J=10\,\zeta(3)-12=0.0206$. This is the supply of combinations we wanted.

\forced The exponent $3$ in $D_n^3$ is forced by the three factors in $\frac1{m^3}$. The kernel
is one natural way to produce those factors symmetrically; Part~II reaches the same combinations
with no integral at all. What we must not do is waste degree: Section~8 shows that even one extra
factor of $D_n$ would be fatal.

\section{Which polynomial? The cancellation requirement}

\q{Question.} Among all integer polynomials of degree $n$, which one makes $J(P)$ small?

\q{Attempt.} Take something simple of degree $n$ that is small on most of the square, say
$P(x)=(1-x)^n$. Its coefficients are integers, the price is $D_n^3$. But
$J\bigl((1-x)^n\bigr)=0.138,\ 0.050,\ 0.017$ at $n=5,10,20$: it decreases like a power of $n$, no
faster than the partial sums did. Multiplied by $D_n^3$ it is $8\cdot10^{8}$ at $n=10$. Worse, for
a carelessly chosen polynomial the sign is not even under control: $P=2-5x$ gives
$J=58\zeta(3)-70=-0.28$, although the kernel is positive.

\q{Lesson.} A polynomial that is merely small does not give exponential smallness. We must use the
freedom in the coefficients to make things \emph{cancel}.

\q{Idea.} Think of one variable at a time: we want $\int_0^1P(x)f(x)\dd x$ small, where $f$ is a
smooth function (the kernel with $y$ fixed, or a piece of it). Smooth functions are well
approximated by polynomials of low degree. So if $P$ gives integral $0$ against \emph{every}
polynomial $g$ of degree less than $n$, then
\[
\int_0^1P\,f=\int_0^1P\,(f-g)\qquad\text{for all such }g,
\]
and the integral only ``sees'' the part of $f$ that polynomials of degree $<n$ cannot reproduce,
which is tiny. This is a design requirement, and it is the most we can ask:

\begin{proposition}[The polynomial is forced]\label{prop:forced}
Fix $n$. Among polynomials of degree at most $n$, those satisfying
$\int_0^1P(x)\,x^j\dd x=0$ for $j=0,1,\dots,n-1$ form a line: there is a non-zero one, it has degree
exactly $n$, and any two are proportional. No non-zero polynomial of degree at most $n$ satisfies one
more condition $\int_0^1P\,x^n=0$.
\end{proposition}
\begin{proof}
These are $n$ linear equations in $n+1$ unknown coefficients, so a non-zero solution exists. If $P$
is a solution and $g$ is any polynomial of degree less than $n$, then $\int Pg=0$. If $P$ itself
had degree less than $n$ we could take $g=P$ and get $\int P^2=0$, so $P=0$. Hence a non-zero
solution has degree exactly $n$; and given two solutions, a combination of them has degree less than
$n$ and is again a solution, so it vanishes. The last sentence is the same argument with $g=P$.
\end{proof}

So \emph{once we decide to spend our $n+1$ coefficients on cancelling everything of lower degree,
there is no further choice}. It remains to find this polynomial and to check that it can be scaled
to have integer coefficients.

\q{Construction.} We need an operation that kills polynomials of degree less than $n$:
differentiating $n$ times does. Integration by parts moves derivatives from one factor to the
other, at the cost of boundary terms. So take a ``seed'' $F$ whose boundary terms all vanish, that
is, $F$ and its first $n-1$ derivatives vanish at both $0$ and $1$. The polynomial of least degree
with this property is $x^n(1-x)^n$. Its $n$-th derivative has degree $n$, as required. Define
\begin{equation}\label{legendre}
P_n(x)=\frac1{n!}\frac{\dd^n}{\dd x^n}\bigl[x^n(1-x)^n\bigr]
=\sum_{k=0}^n(-1)^k\binom nk\binom{n+k}k x^k ,
\end{equation}
where the formula comes from expanding $(1-x)^n$ and differentiating term by term. The division by
$n!$ is exactly what makes the coefficients integers with no common factor ($P_n(0)=1$):
\[
P_1=1-2x,\qquad P_2=1-6x+6x^2,\qquad P_3=1-12x+30x^2-20x^3 .
\]
These are called shifted Legendre polynomials; no knowledge of them is assumed. Integrating by parts
$n$ times, with no boundary terms, gives the formula on which everything rests: for any $f$ with
$n$ continuous derivatives on $[0,1]$,
\begin{equation}\label{ibp}
\int_0^1P_n(x)f(x)\dd x=\frac{(-1)^n}{n!}\int_0^1x^n(1-x)^n\,f^{(n)}(x)\dd x .
\end{equation}
If $f$ is a polynomial of degree less than $n$, then $f^{(n)}=0$: so $P_n$ is the polynomial of
Proposition~\ref{prop:forced}.

\q{Why not use the seed itself?} The seed $x^n(1-x)^n$ is positive and at most $4^{-n}$, so
$J$ of it is positive and at most $2\zeta(3)\,16^{-n}$. That looks like everything we want. But
its degree is $2n$, so the price is $D_{2n}^3$, about $e^{6n}\approx400^n$. At $n=10$:
$J=1.3\cdot10^{-13}$, $D_{20}^3=1.3\cdot10^{25}$, product $10^{12}$. In general it cannot work: the
binomial coefficient $\binom{2n}n\ge\frac{4^n}{2n+1}$ divides $D_{2n}$ (a classical fact about binomial coefficients), so $D_{2n}^3$ is at least
about $64^n$, while $J$ of the seed is not smaller than about $16^{-n}$ (the seed is close to
$4^{-n}$ near $x=\frac12$). Differentiating $n$ times is what brings the degree, and the price,
back down, while formula \eqref{ibp} keeps the seed available whenever we need it.

\q{The numbers we shall study.} From now on
\[
I_n=J(P_n)=2\bigl(a_n\zeta(3)-b_n\bigr),\qquad
a_n=\sum_{k=0}^n\binom nk^2\binom{n+k}k^2,\qquad D_n^3\,b_n\in\Z,
\]
by Proposition~\ref{prop:arith}: $a_n$ is the sum of the squares of the coefficients of $P_n$, and
$b_n=\frac12\beta(P_n)$. The first values are $a_n=1,5,73,1445$ and
$b_n=0,\,6,\,\frac{351}4,\,\frac{62531}{36}$:
\[
I_0=2\zeta(3),\quad I_1=10\zeta(3)-12=0.0206,\quad I_2=146\zeta(3)-\tfrac{351}2=0.00031,\quad
I_3=5.7\cdot10^{-6}.
\]
Compare $n=10$: the careless $(1-x)^{10}$ gave $0.05$; $P_{10}$ gives $I_{10}=2.2\cdot10^{-17}$,
at the same price. The integers $a_n$ (the Ap\'ery numbers) were not guessed: they are the
diagonal of $P_n(x)P_n(y)$, and only the diagonal carries $\zeta(3)$.

\forced The decision to cancel all lower degrees is a choice (it is the natural way to use
$n+1$ coefficients, and it is what makes the next section possible). Given that decision, $P_n$
is forced up to sign by Proposition~\ref{prop:forced} and the requirement of integer coefficients
without common factor; and then $a_n$, $b_n$, $I_n$ are forced too.

\section{Why is $I_n$ not zero? Making the sign visible}

\q{Question.} The numbers above suggest that $I_n$ is small and positive. But $P_n(x)P_n(y)$
changes sign all over the square, and we saw with $2-5x$ that positivity of the kernel proves
nothing. How can we see the sign of $I_n$?

\q{Idea.} Formula \eqref{ibp} trades the oscillating $P_n$ for the positive seed $x^n(1-x)^n$,
\emph{provided} we can differentiate the other factor $n$ times and control its sign. For one very
simple function this is easy: if $f(x)=\frac1{x+t}$ with $t>0$, then
$f^{(n)}(x)=\frac{(-1)^nn!}{(x+t)^{n+1}}$, and the signs and factorials in \eqref{ibp} cancel
exactly:
\begin{equation}\label{onevar}
\int_0^1\frac{P_n(x)}{x+t}\dd x=\int_0^1\frac{x^n(1-x)^n}{(x+t)^{n+1}}\dd x>0 .
\end{equation}
(With $t=1$ this already proves a theorem: the left side is an integer combination of $1$ and
$\log2$ with denominator $D_n$, the right side is positive and tiny, and $\log2$ is irrational.)

\q{Obstacle.} Our kernel $\frac{-\log(xy)}{1-xy}$ is not of this simple form, and its $n$-th
derivatives are a mess with no visible sign.

\q{Construction.} Write the kernel as a \emph{mixture of products of simple fractions}. A logarithm
is an integral of a reciprocal: for $a,b>0$, splitting into partial fractions,
$\int_0^\infty\frac{\dd t}{(t+a)(t+b)}=\frac{\log b-\log a}{b-a}$. We want the numerator
$-\log(xy)$ and the denominator $1-xy$; taking $a=x$ and $b=\frac1y$ gives both, and
\begin{equation}\label{sep}
K(x,y)=\int_0^\infty\frac{\dd t}{(x+t)(1+yt)}\qquad(0<x,y<1).
\end{equation}
For each fixed $t$ the variables are now separated, each sitting in a simple fraction.

\q{Carrying it out.} Insert \eqref{sep} into $I_n=\iint P_n(x)P_n(y)K$. For fixed $t$ the integral
over $x$ and $y$ is a product of two one-variable integrals. The first is \eqref{onevar}. The
second is the same computation with $f(y)=\frac1{1+yt}$, whose $n$-th derivative is
$\frac{(-1)^nn!\,t^n}{(1+yt)^{n+1}}$. Therefore
\begin{equation}\label{triple}
\boxed{\;I_n=\int_0^\infty\!\!\int_0^1\!\!\int_0^1
\frac{\bigl[x(1-x)\,y(1-y)\,t\bigr]^n}{\bigl[(x+t)(1+yt)\bigr]^{n+1}}\dd x\dd y\dd t\ >\ 0 .\;}
\end{equation}
Every factor is positive. So $I_n>0$ for every $n$: the cancellation has been converted into a
positive integral \emph{before} any estimate is made. (The order of integration may be changed
because $|P_n(x)P_n(y)|/((x+t)(1+yt))$ has total integral at most $\max|P_n|^2\cdot2\zeta(3)$.)

\forced The extra variable $t$ is the price of making the kernel differentiable in closed form.
The substitution $t=\frac{1-z}{yz}$ turns \eqref{triple} into Beukers' original integral over the
cube, $\int_{[0,1]^3}\frac{[x(1-x)y(1-y)z(1-z)]^n}{(1-z+xyz)^{n+1}}$; the two are the same.

\section{How small is $I_n$? Where $\sqrt2$ comes from}

\q{Question.} The integrand in \eqref{triple} is an $n$-th power. How fast does it shrink?

Write it as $G^n\cdot\frac1Q$ with
\[
Q=(x+t)(1+yt),\qquad G=\frac{x(1-x)\,y(1-y)\,t}{Q}.
\]
If $G\le\rho$ everywhere, then $I_n\le\rho^n\iiint\frac1Q=\rho^n\iint K=2\zeta(3)\rho^n$. So the
whole question is: \emph{what is the maximum of $G$?} It is a competition: the numerator wants $x$
and $y$ near $\frac12$, the denominator wants them near $0$.

\emph{Reduce to one variable.} Everything is symmetric, so the maximum should have $x=y$ and a
balanced $t$. Put $s=\sqrt{xy}$. By AM--GM, $x+y\ge2s$, so
$(1-x)(1-y)=1-(x+y)+xy\le(1-s)^2$; and $\frac Qt=1+xy+\frac xt+yt\ge1+s^2+2s=(1+s)^2$. Hence
\[
G\le\Bigl(\frac{s(1-s)}{1+s}\Bigr)^2 .
\]
\emph{One-variable maximum.} The function $\frac{s(1-s)}{1+s}$ has derivative zero where
$s^2+2s-1=0$, that is, at $s=\sqrt2-1$. This is where $\sqrt2$ enters the subject. Its value there
is $(\sqrt2-1)^2$, as the identity $\alpha^2(1+s)-s(1-s)=(s-\alpha)^2\ge0$ for $\alpha=\sqrt2-1$
confirms. Therefore
\begin{equation}\label{boundA}
\boxed{\;0<I_n\le2\zeta(3)\,\rho^{\,n},\qquad\rho=(\sqrt2-1)^4=17-12\sqrt2=0.0294\ldots\;}
\end{equation}
So $I_n$ shrinks by a factor of about $34$ at each step. Since $G=\rho$ is attained (at
$x=y=\sqrt2-1$, $t=1$), no better rate is possible for these numbers.

\section{Is the price affordable? The size of $D_n$}

\q{Question.} We must compare $D_n^3$ with $\rho^{-n}=33.97^n$. We need $D_n<C^n$ with
$C^3<33.97$, that is $C<3.238$.

\q{What is true.} $D_n$ grows like $e^n=2.718^n$ (for example $D_{12}^{1/12}=2.35$). That
statement is the prime number theorem, which is deep. We need far less, for a reason worth
noticing: since $I_n>0$ for \emph{every} $n$, it is enough to have a good bound on $D_n$ for
\emph{infinitely many} $n$. And that follows from one exact identity.

\begin{lemma}\label{lem:fact}
For every positive integer $N$,\quad $N!=D_N\cdot D_{\lfloor N/2\rfloor}\cdot D_{\lfloor N/3\rfloor}
\cdots D_{\lfloor N/N\rfloor}$.
\end{lemma}
\noindent Example: $6!=D_6D_3D_2D_1D_1D_1=60\cdot6\cdot2=720$.
\begin{proof}
Fix a prime $p$ and count how many times it divides each side. $D_m$ contains $p$ once for each
power $p^a\le m$. So the right side contains $p$ once for each pair $(k,a)$ with
$p^a\le\lfloor N/k\rfloor$, that is, with $k\le N/p^a$. For each $a$ there are $\lfloor N/p^a\rfloor$
such $k$. And $\sum_a\lfloor N/p^a\rfloor$ is how many times $p$ divides $N!$ (count the multiples of
$p$, then of $p^2$, and so on).
\end{proof}

\begin{lemma}\label{lem:lcm}
For every $C>e$ there are infinitely many $n$ with $D_n<C^n$. In particular $D_n<3^n$ for
infinitely many $n$.
\end{lemma}
\noindent\emph{The idea.} By Lemma~\ref{lem:fact}, $N!$ is a product of the $D$'s at the scales
$N,\frac N2,\frac N3,\dots$ If all of them were at least $C^{\text{(index)}}$, the product would be
at least $C$ to the power $N+\frac N2+\frac N3+\cdots\approx N\log N$. But $N!\le N^N=e^{N\log N}$. So
$C\le e$. The $D$'s cannot all be above the average that the factorial imposes.
\begin{proof}
Suppose the lemma fails for some $C>e$: then $D_m\ge C^m$ for all $m\ge M$, for some $M$. Take $N$
large and let $L=\lfloor N/M\rfloor$. For $k\le L$ the index $\lfloor N/k\rfloor$ is at least $M$, so
\[
N\log N\ \ge\ \log N!\ \ge\ \log C\sum_{k=1}^{L}\Bigl\lfloor\frac Nk\Bigr\rfloor
\ \ge\ \log C\,\Bigl(N\bigl(1+\tfrac12+\cdots+\tfrac1L\bigr)-L\Bigr).
\]
Since $1+\frac12+\cdots+\frac1L\ge\log(L+1)>\log N-\log M$ and $L\le N/M$, dividing by $N$ gives
$\log N\ge\log C\,(\log N-\log M-\frac1M)$. As $\log C>1$ and $M$ is fixed, this is false for large
$N$.
\end{proof}

\section{The proof, and how little room there is}

\begin{theorem}[Ap\'ery]
$\zeta(3)$ is irrational.
\end{theorem}
\begin{proof}
Suppose $\zeta(3)=p/q$. By Proposition~\ref{prop:arith},
$q\,D_n^3I_n=2\bigl(D_n^3a_n\,p-D_n^3b_n\,q\bigr)$ is an integer. By \eqref{triple} it is positive.
By \eqref{boundA} and Lemma~\ref{lem:lcm}, for infinitely many $n$ it is less than
$2q\zeta(3)\,(27\rho)^n$. Since $\rho=\frac1{17+12\sqrt2}<\frac1{32}$ (because $\sqrt2>\frac54$),
this is less than $2q\zeta(3)(\frac{27}{32})^n$, which is below $1$ for large $n$. A positive
integer below $1$ does not exist.
\end{proof}

\noindent For orientation, here are the actual values of $D_n^3I_n$:
\begin{center}\small
\begin{tabular}{@{}lcccccc@{}}
\toprule
$n$ & 1 & 2 & 4 & 6 & 8 & 10\\
$D_n^3\,I_n$ & $0.021$ & $0.0025$ & $2.0\cdot10^{-4}$ & $1.3\cdot10^{-5}$ & $2.1\cdot10^{-5}$ &
$3.6\cdot10^{-7}$\\
\bottomrule
\end{tabular}
\end{center}
They are not monotone ($D_n$ jumps when $n$ is a prime power) but they decay.

\q{What did what.}
\begin{center}\small
\begin{tabular}{@{}p{43mm}p{52mm}p{52mm}@{}}
\toprule
object & the need it answers & forced or chosen\\
\midrule
criterion (Lemma~\ref{lem:crit}) & a rigid property of fractions & forced: equivalent to
irrationality\\
$D_n$ instead of $n!$ & cheapest common denominator & forced\\
kernel $\frac{-\log(xy)}{1-xy}$ & three factors $\frac1{m+1}$; only $\zeta(3)$ appears &
natural; Part~II avoids it\\
degree $n$ in each variable & price $D_n^3$ and no more & forced: degree $2n$ fails
(Section~4)\\
$P_n$ & cancel everything of lower degree & a choice of principle; then forced
(Proposition~\ref{prop:forced})\\
$a_n=\sum\binom nk^2\binom{n+k}k^2$ & nothing: it is the diagonal of $P_n(x)P_n(y)$ & forced by
$P_n$\\
identity \eqref{sep} & a kernel that can be differentiated $n$ times & one device; Part~II has
others\\
$(\sqrt2-1)^4$ & maximum of the integrand & forced: it is the exact rate\\
subsequence of $n$ & avoid the prime number theorem & available because $I_n>0$ for all $n$\\
\bottomrule
\end{tabular}
\end{center}

\q{How little room there is.} The proof closes because $e^3=20.1<33.97$. Change any one thing and
it does not:
\begin{itemize}
\item one more factor of $D_n$ (price $D_n^4\approx54.6^n$): fails, since $54.6>33.97$;
\item degree $2n$ instead of $n$: fails (Section~4);
\item $n!$ instead of $D_n$: fails at once, $n!$ grows faster than any exponential;
\item the cube of the logarithm in the kernel, which would give $\zeta(5)$ with denominators
$D_n^5$: the cancellation is then lost, and the combinations are only of size $n^{-4}\log n$
(reported in the companion document).
\end{itemize}
This is why $\zeta(3)$ is special for this method, and why every economy above was necessary.

\part*{Part II. The same objects without integrals}

Zudilin's paper proves the theorem starting from a rational function,
$\bigl(\frac{(t-1)\cdots(t-n)}{t(t+1)\cdots(t+n)}\bigr)^2$, with no explanation of where it comes
from. Part~II shows that it is $P_n$ in disguise, so that nothing new has to be guessed; and that
this second language explains the famous recurrence and gives a second proof of positivity and
smallness.

\section{The polynomial $P_n$ seen as a rational function}

\q{Question.} The series $\zeta(3)=\sum\frac1{t^3}$ is a sum over integers $t$. Can we describe
$I_n$ as a sum over $t$ as well?

\q{Idea.} Look at all the moments of $P_n$ at once: integrate it against $x^{t-1}$, where the
exponent $t>0$ is a variable. Term by term, from \eqref{legendre},
\[
\int_0^1x^{t-1}P_n(x)\dd x=\sum_{k=0}^n\frac{(-1)^k\binom nk\binom{n+k}k}{t+k}.
\]
What kind of function of $t$ is this? Three facts determine it completely:
\begin{itemize}
\item it is a rational function with simple poles at $t=0,-1,\dots,-n$ (because $P_n$ has degree
$n$) and numerator of degree at most $n$;
\item it vanishes at $t=1,2,\dots,n$ (because $P_n$ is orthogonal to $1,x,\dots,x^{n-1}$);
\item its residue at $t=0$ is $P_n(0)=1$.
\end{itemize}
A numerator of degree at most $n$ with $n$ prescribed zeros is determined up to a constant, and
the third fact fixes the constant. Hence
\begin{equation}\label{mellin}
\int_0^1x^{t-1}P_n(x)\dd x=(-1)^n r_n(t),\qquad
r_n(t)=\frac{(t-1)(t-2)\cdots(t-n)}{t(t+1)\cdots(t+n)} .
\end{equation}
\emph{The zeros of $r_n$ are the orthogonality of $P_n$; its poles are the degree of $P_n$.}
Comparing the two displays also gives the partial fractions of $r_n$:
\begin{equation}\label{pf}
r_n(t)=\sum_{k=0}^n\frac{c_k}{t+k},\qquad c_k=(-1)^{n-k}\binom nk\binom{n+k}k\in\Z .
\end{equation}

\q{Why the square, and why a derivative?} Since $I_n$ contains $P_n(x)P_n(y)$, the transform
appears twice: $\iint(xy)^{t-1}P_n(x)P_n(y)=r_n(t)^2$. Differentiating $(xy)^{t-1}$ with respect to
$t$ produces $\log(xy)$, and summing $(xy)^{t-1}$ over $t\ge1$ produces $\frac1{1-xy}$. So, with
$R_n(t)=r_n(t)^2$,
\begin{equation}\label{defI}
I_n=-\sum_{t=1}^\infty R_n'(t).
\end{equation}
(The interchanges are justified because $P_n$ is bounded and $K$ is positive and integrable.) This
is exactly Zudilin's starting formula. From the point of view of sums one can also see directly
why this combination is the right one:
\begin{itemize}
\item we need poles of order $3$ to reach $\sum\frac1{t^3}$: squaring $r_n$ gives order $2$, and
one differentiation gives order $3$;
\item squaring makes the zeros at $t=1,\dots,n$ double, so they survive the differentiation: the
first $n$ terms of \eqref{defI} vanish;
\item the alternative with poles of order $3$, the cube $r_n^3$, fails: its sum also contains
$\zeta(2)$, with a coefficient that does not vanish (it is $18$ for $n=1$), so it would be a
combination of two unknown constants. With the derivative of a square, $\zeta(2)$ cancels by itself,
as we now check.
\end{itemize}

\q{The arithmetic again, in this language.} Squaring \eqref{pf} and using
$\frac1{(t+k)(t+j)}=\frac1{j-k}\bigl(\frac1{t+k}-\frac1{t+j}\bigr)$,
\[
R_n(t)=\sum_{k=0}^n\Bigl[\frac{c_k^2}{(t+k)^2}+\frac{e_k}{t+k}\Bigr],\qquad
e_k=2c_k\sum_{j\ne k}\frac{c_j}{j-k}.
\]
Then $-R_n'(t)=\sum_k\bigl[\frac{2c_k^2}{(t+k)^3}+\frac{e_k}{(t+k)^2}\bigr]$, and summing over
$t\ge1$ with $\sum_{t\ge1}(t+k)^{-j}=\zeta(j)-H_k^{(j)}$:
\[
I_n=2\Bigl(\sum_kc_k^2\Bigr)\zeta(3)+\Bigl(\sum_ke_k\Bigr)\zeta(2)
-\sum_k\bigl[2c_k^2H_k^{(3)}+e_kH_k^{(2)}\bigr].
\]
Now $\sum_ke_k=0$, because in the double sum the terms $(j,k)$ and $(k,j)$ cancel. So $\zeta(2)$
disappears, $a_n=\sum c_k^2$ as before, and since $D_ne_k$ is an integer the rational part has
denominator dividing $D_n^3$. This re-proves Proposition~\ref{prop:arith} for $P_n$, and gives the
explicit formula $b_n=\sum_k\bigl[c_k^2H_k^{(3)}+\tfrac12e_kH_k^{(2)}\bigr]$.

\section{The recurrence: how to find it, and what it gives}

\q{Question.} Ap\'ery's original proof rested on the fact that $a_n$ and $b_n$ both satisfy
\[
(n+1)^3u_{n+1}-(2n+1)(17n^2+17n+5)\,u_n+n^3u_{n-1}=0 .
\]
Where could such a thing come from, and what is it good for?

\q{What it is good for.} Two different solutions of one second-order recurrence are tied together
by a simple identity (below), and that identity will give the sign and the size of $I_n$ with no
integral at all. So it is a second, independent route to Sections~5 and~6.

\q{How to find it.} Both $a_n$ and $I_n$ are obtained from $R_n$ by an operation that kills any
function of the form $S(t+1)-S(t)$: for $I_n$ the operation is ``differentiate and sum over
$t\ge1$'', which telescopes. So we look for constants $A,B,C$ and a function $S$ with
\[
A\,R_{n+1}(t)+B\,R_n(t)+C\,R_{n-1}(t)=S(t+1)-S(t).
\]
The simplest guess is $S=s\cdot R_n$ with $s$ a polynomial. Its degree can be at most $2$: since
$R_n(t)\approx t^{-2}$, a polynomial of degree $2$ makes $sR_n$ tend to a constant, which the
difference kills, while degree $3$ would leave something that does not tend to $0$. Dividing by
$R_n(t)$, all the ratios are explicit,
\[
\frac{R_{n+1}(t)}{R_n(t)}=\Bigl(\frac{t-n-1}{t+n+1}\Bigr)^2,\quad
\frac{R_{n-1}(t)}{R_n(t)}=\Bigl(\frac{t+n}{t-n}\Bigr)^2,\quad
\frac{R_n(t+1)}{R_n(t)}=\frac{t^4}{(t-n)^2(t+n+1)^2},
\]
and after multiplying by $(t-n)^2(t+n+1)^2$ the requirement becomes an identity between polynomials
of degree $4$ in $t$:
\begin{multline}\label{poly}
A(t-n-1)^2(t-n)^2+B(t-n)^2(t+n+1)^2+C(t+n)^2(t+n+1)^2\\
=s(t+1)\,t^4-s(t)\,(t-n)^2(t+n+1)^2 .
\end{multline}
(On the right the terms of degree $6$ and $5$ cancel automatically.) That is $5$ linear equations,
one for each coefficient, in $6$ unknowns: $A,B,C$ and the three coefficients of $s$. We solved this
system with $n$ as a symbol: it has rank $5$, so \emph{the solution is unique up to a common
factor}. It is
\begin{gather*}
A=(n+1)^3,\qquad B=-(2n+1)(17n^2+17n+5),\qquad C=n^3,\\
s_n(t)=4(2n+1)\bigl(-2t^2+t+(2n+1)^2\bigr).
\end{gather*}
So the cubes and the number $17$ are not guesses and not a miracle: they are the unique solution of
a small linear system that anyone can write down.

\begin{proposition}\label{prop:cert}
With these values, for $n\ge1$,
\begin{equation}\label{cert}
(n+1)^3R_{n+1}(t)-\beta_nR_n(t)+n^3R_{n-1}(t)=s_n(t+1)R_n(t+1)-s_n(t)R_n(t),
\end{equation}
where $\beta_n=(2n+1)(17n^2+17n+5)$.
\end{proposition}
\begin{proof}
We must verify \eqref{poly}. Both sides have degree at most $4$, so it is enough to find five zeros
of the difference, counted with multiplicity. Write $\mu=2n+1$ and $u(t)=-2t^2+t+\mu^2$, so
$s_n=4\mu u$; note $u(n+1)=n\mu$, $u(-n)=(n+1)\mu$, $u'(n+1)=-(4n+3)$, $u'(-n)=4n+1$.
\begin{itemize}
\item $t=n$: left $=n^3(2n)^2\mu^2=4n^5\mu^2$; right $=4\mu\,u(n+1)\,n^4=4n^5\mu^2$.
\item $t=-n-1$: left $=(n+1)^3(2n+2)^2\mu^2$; right $=4\mu\,u(-n)(n+1)^4$; both are
$4(n+1)^5\mu^2$.
\item $t=0$: left $=n^2(n+1)^2[(n+1)^3+n^3-\beta_n]$, right $=-4\mu^3n^2(n+1)^2$. They agree because
$(n+1)^3+n^3=\mu(n^2+n+1)$ and $17n^2+17n+5-(n^2+n+1)=4\mu^2$.
\item derivative at $t=n$ (only terms without the factor $(t-n)^2$ contribute): left
$=n^3\cdot2(2n)\mu(4n+1)$; right $=4\mu[u'(n+1)n^4+4u(n+1)n^3]$; both are $4n^4\mu(4n+1)$.
\item derivative at $t=-n-1$: left $=(n+1)^3\cdot2(-2n-2)(-\mu)(-4n-3)$; right
$=4\mu[u'(-n)(n+1)^4-4u(-n)(n+1)^3]$; both are $-4(n+1)^4\mu(4n+3)$.\qedhere
\end{itemize}
\end{proof}

\q{From the identity to two recurrences.} Write
$\mathcal L[X]_n=(n+1)^3X_{n+1}-\beta_nX_n+n^3X_{n-1}$.
\begin{itemize}
\item \emph{$\mathcal L[I]_n=0$ for $n\ge1$.} Differentiate \eqref{cert} and sum over
$t=1,\dots,T$. The right side telescopes to $(s_nR_n)'(T+1)-(s_nR_n)'(1)$. The last term is $0$
because $R_n$ has a double zero at $t=1$. The other tends to $0$ because $s_nR_n$ tends to a
constant, so its derivative is $O(T^{-2})$.
\item \emph{$\mathcal L[a]_n=0$ for $n\ge1$.} Recall $a_n=\sum c_k^2$ is the sum of the
coefficients of the double poles of $R_n$. For any rational function $f$, let $\delta(f)$ be the
sum of all its double-pole coefficients (in its partial fractions, which are unique). This
operation is linear and does not notice a shift $t\mapsto t+1$, which only moves the poles. So
$\delta$ of the right side of \eqref{cert} is $0$, and $\delta(R_n)=a_n$.
\end{itemize}
(We do not get the recurrence for $a_n$ by ``reading off the coefficient of $\zeta(3)$'' in the
recurrence for $I_n$: that would assume $1$ and $\zeta(3)$ are independent over the rationals,
which is what we are proving.)

\q{The identity that ties two solutions together.} For a differential equation $y''=\dots$, two
solutions $y_1,y_2$ have a Wronskian $y_1y_2'-y_1'y_2$ that obeys a first-order equation. The same
happens for recurrences. Put $\Omega_n=a_nI_{n-1}-a_{n-1}I_n$. Using both recurrences,
\[
(n+1)^3\Omega_{n+1}=(\beta_na_n-n^3a_{n-1})I_n-a_n(\beta_nI_n-n^3I_{n-1})=n^3\,\Omega_n .
\]
So $n^3\Omega_n$ is constant; and $\Omega_1=a_1I_0-a_0I_1=5\cdot2\zeta(3)-(10\zeta(3)-12)=12$.
Dividing by $a_{n-1}a_n$:
\begin{equation}\label{cas}
\frac{I_{n-1}}{a_{n-1}}-\frac{I_n}{a_n}=\frac{12}{n^3\,a_{n-1}a_n}>0\qquad(n\ge1).
\end{equation}
\emph{The sequence $I_n/a_n$ decreases, by an amount we know exactly.} If it also tends to $0$, then
every term is positive and equals the sum of all later decrements.

\q{Second proof that $I_n$ is positive and small.}
\begin{enumerate}
\item \emph{$I_n/a_n\to0$.} A crude bound suffices. For integers $t>n\ge1$: $0<r_n(t)<\frac1t$
(each factor $\frac{t-j}{t+j}$ is less than $1$), so $R_n(t)<t^{-2}$; and $R_n'=2R_n\varphi_n$ with
$\varphi_n(t)=\sum_{j=1}^n\frac1{t-j}-\sum_{j=0}^n\frac1{t+j}$, a difference of two positive numbers
each at most $1+\frac12+\dots+\frac1n\le n$. As the first $n$ terms of \eqref{defI} vanish,
$|I_n|\le2n\sum_{t>n}t^{-2}<2$. And $a_n\to\infty$, since it contains the term $n^2(n+1)^2$.
\item \emph{Sign.} Summing \eqref{cas} from $n+1$ to infinity:
$\displaystyle\frac{I_n}{a_n}=\sum_{m>n}\frac{12}{m^3a_{m-1}a_m}>0$.
\item \emph{Size.} The $a_n$ increase with $n$ (each term $\binom nk^2\binom{n+k}k^2$ does not
decrease). So $a_{m-1}\ge a_n$ and $a_m\ge a_{n+1}$ for $m>n$, and
$I_n\le\frac{12}{a_{n+1}}\sum_{m>n}\frac1{m^3}<\frac6{n^2a_{n+1}}$. \emph{All the smallness of $I_n$
is the largeness of the integer $a_n$.}
\item \emph{Growth of $a_n$.} From the recurrence and $a_{n-1}\le a_n$,
$\frac{a_{n+1}}{a_n}\ge\frac{\beta_n-n^3}{(n+1)^3}=33-\frac{48n^2+72n+28}{(n+1)^3}>33-\frac{48}{n+1}
\ge28$ for $n\ge9$. So $a_n\ge c_0\,28^n$ for some $c_0>0$, and $0<I_n<\frac6{c_0}28^{-n}$.
\end{enumerate}
Since $28>27$, the proof of the Theorem in Section~8 goes through word for word with this in place
of \eqref{triple} and \eqref{boundA}. Taking $n=0$ in step~2 gives a series for the constant itself:
$\zeta(3)=6\sum_{m\ge1}\frac1{m^3a_{m-1}a_m}$.

\begin{remark}[What Zudilin's paper does instead]
Zudilin proves the arithmetic and identity \eqref{cert} as above. For positivity and size he
introduces a second series (due to Ball) whose terms are visibly positive, bounds it term by term,
and then proves it equals $I_n$ by showing that it satisfies the same recurrence, which needs a
second identity of the type \eqref{cert} with a much longer polynomial. Steps 1--4 above replace
all of that. The price is the rate $28^{-n}$ instead of the true $33.97^{-n}$, which is still
enough.
\end{remark}

\section{A third view: why the same integer measures arithmetic and size}

Step~3 above says $I_n\,a_n$ is only polynomially small. There is a direct reason, which needs
neither the triple integral nor the recurrence. We include it because it explains the role of
$a_n$; it can be skipped.

\q{Idea.} $P_n$ is characterised by a differential equation: with $Jf=-(x(1-x)f')'$ one has
$JP_n=n(n+1)P_n$. (Reason: $J$ does not raise degrees and $\int(Jf)g=\int f(Jg)$, so $JP_n$ is again
orthogonal to all lower degrees, hence a multiple of $P_n$.) The kernel, as a function $K(v)$ of
$v=xy$, also satisfies a differential equation. Together they give an equation for the function
\[
k_n(x)=\int_0^1P_n(y)\,K(xy)\dd y,\qquad\text{for which}\quad I_n=\int_0^1P_n(x)k_n(x)\dd x .
\]
Write $\lambda_n=n(n+1)$ and $\mathcal Hf=-(x^2(1-x)f')'+xf$. Three direct computations:
(i) $(v^2-v)K''+(3v-1)K'+K=0$, which implies $\mathcal H_xK(xy)=-J_yK(xy)$;
(ii) $\mathcal Hx^j=(j+1)^2x^{j+1}-j(j+1)x^j$;
(iii) $vK'(v)\to-1$ as $v\to0$.

\q{Step 1: an equation for $k_n$.} By (i), $\mathcal Hk_n=-\int_0^1P_n(y)J_yK(xy)\dd y$. Move $J_y$
onto $P_n$ by two integrations by parts. This gives $\lambda_nk_n$ plus boundary terms; the only one
that survives is at $y=0$, where $P_n(0)=1$ and $y\,\partial_yK(xy)\to-1$ by (iii). The result is
\begin{equation}\label{kn}
(\mathcal H+\lambda_n)\,k_n=1 .
\end{equation}

\q{Step 2: where $a_n$ comes from.} By (ii), $\mathcal H+\lambda_n$ raises the degree of a
polynomial by exactly one. So $P_n$ can be written uniquely as $(\mathcal H+\lambda_n)w+\kappa$ with
$w$ a polynomial of degree less than $n$ and $\kappa$ a constant. To find $\kappa$, apply the linear
functional $E(x^j)=(-1)^j\binom nj\binom{n+j}j$: by (ii) and the ratio
$\binom n{j+1}\binom{n+j+1}{j+1}\big/\binom nj\binom{n+j}j=\frac{(n-j)(n+j+1)}{(j+1)^2}$ it kills
$(\mathcal H+\lambda_n)x^j$ for $j<n$, and $E(1)=1$. Hence
$\kappa=E(P_n)=\sum_j\binom nj^2\binom{n+j}j^2=a_n$: the coefficients of $P_n$ are paired with
themselves.

\q{Step 3: an energy identity.} Let $g=w+a_nk_n$. Then $(\mathcal H+\lambda_n)g=P_n$ by
\eqref{kn}, and $\int P_ng=a_nI_n$ because $w$ has degree less than $n$. Multiply
$(\mathcal H+\lambda_n)g=P_n$ by $g$ and integrate by parts once (the boundary terms vanish):
\[
a_nI_n=\int_0^1x^2(1-x)\,g'^2\dd x+\int_0^1(x+\lambda_n)\,g^2\dd x\ >\ 0 .
\]
So $I_n>0$. And by Cauchy--Schwarz with $\int_0^1P_n^2=\frac1{2n+1}$:
$\lambda_n\!\int g^2\le a_nI_n=\int P_ng\le\bigl(\frac1{2n+1}\int g^2\bigr)^{1/2}$, which gives
\[
\boxed{\;0<I_n\le\frac1{n(n+1)(2n+1)\,a_n}\qquad(n\ge1).\;}
\]
With step~4 of Section~10 (or a direct lower bound for $a_n$) this is a third complete proof of
(B). The full details of the boundary terms are in the Codex kernel-energy note cited below.

\section{Status, sources, checks}

\paragraph{What is proved here.} Part~I is a complete proof. Section~10 is a second complete
proof of positivity and smallness. Section~11 is complete except that the vanishing of two boundary
terms is asserted and not written out.

\paragraph{Sources.} Part~I follows the Codex manuscript in
\texttt{Zeta3-4-Motivated-proof-Codex/}, which is a motivated account of Beukers' proof; the
comparison numbers in Sections~2, 4 and~8 were computed for this document. Section~9 is the
dictionary found by both agents. Section~10 is Claude's simplification of Zudilin's proof
(\texttt{Zeta3-2-Elementary-proof/}), with the search for \eqref{cert} as in the Codex note
\texttt{simplification.md}. Section~11 is from \texttt{Codex-Zeta3-03-Kernel-energy/}. Of the
original papers only Zudilin's (arXiv:math/0202159) was retrieved and read; attributions to Ap\'ery,
Beukers and Ball are as given there.

\paragraph{Checks (AI only).} The mathematical content of Sections 3--11 is that of the previous
version of this document, in which a second reader redid every formula independently and found no
error. For this rewrite the following were computed afresh: all numerical values quoted; the
moments and the forms $J(P)$ exactly; the rank of the linear system behind \eqref{poly} with $n$
symbolic (rank $5$, solution as stated); the coefficient $18$ of $\zeta(2)$ for $r_1^3$. The new
text was then read by a fresh reader for correctness and for whether each object is motivated
before it is used. \textbf{No human has read this document.}

\end{document}
